\documentclass[11pt]{ctexart}
\usepackage[a4paper,margin=1in]{geometry}
\usepackage{graphicx}
\usepackage{xcolor}
\usepackage{hyperref}
\definecolor{refgray}{gray}{0.45}
\title{\textbf{市场最新观点汇总}\\[0.4em]{\large 通胀预期重夺定价权，AI叙事与宏观风险激烈碰撞，供给侧约束成为多板块核心矛盾}}
\author{KC桌面 自动生成}
\date{260520}
\begin{document}
\maketitle
\section*{一页摘要}
\begin{itemize}
  \item 通胀担忧取代AI主题成为短期市场核心矛盾，能源板块上涨而利率敏感板块承压，新兴市场资金外流压力加大。
  \item 美债多头全线撤退，空头主导市场，资金主动缩短久期暴露，利率定价逻辑从交易降息预期转向交易供给和仓位结构。
  \item 亚洲正经历由AI基建、能源转型、国防开支和工业资本开支共同驱动的资本支出超级周期，工业周期超越能源冲击成为增长主导变量。
  \item 铜冶炼加工费跌至-104美元/吨，铝市场遭遇近50年最大供给冲击之一，供给侧不可逆重置正在重塑有色金属定价逻辑。
  \item 中国房地产广义库存已下降约19\%，供给端结构性收缩为更可持续的复苏奠定基础，二手房市场活跃度远超新房。
  \item AI硬件供应链从需求侧转向供给侧，产能瓶颈、组件短缺和多源采购策略加速，光学引擎渗透率跃迁和CPU重新定价成为新焦点。
  \item 机构资金从押注最大市值AI赢家转向重仓基础设施瓶颈环节，存储与半导体设备成为超配方向。
\end{itemize}
\section{宏观与利率}
\textbf{通胀预期重新成为资产定价主导变量，利率市场多头叙事周期终结，美联储缩表影响被高估。}

\begin{itemize}
  \item 全球股市上周下跌0.6\%，能源板块上涨4.2\%而公用事业、房地产、材料板块跌幅均超2.9\%，通胀担忧正在取代AI叙事成为短期市场核心矛盾。
  \item 新兴市场跌幅最大（-2.5\%），对全球利率预期高度敏感，若通胀数据进一步超预期可能面临更大资金外流压力。
  \item 美债期货市场上虚值多头仓位被大规模平仓，空头占据主导且大部分处于实值状态，利率定价逻辑已从交易降息预期切换至交易供给和仓位结构。
  \item 资金流入集中于短期政府债和投资级信用债而非长久期国债，投资者主动缩短久期暴露以应对可能继续加息的环境。
  \item 美联储新主席Warsh对缩表的实际影响可能雷声大雨点小，负债端三大项中仅准备金有操作空间，最可能走亲银行路线通过放松监管减少准备金需求。
  \item 全球盈利周期仍在改善，新闻情绪面也在好转，但通胀上行可能迫使央行重新加息，当前市场定价中隐含的软着陆叙事可能低估了这一尾部风险。
\end{itemize}
\begin{figure}[htbp]
\centering
\includegraphics[width=0.88\linewidth]{figures/fig_005.jpg}
\caption{Exhibit 1 - Exhibit 1: Curve-o-meter Rates to trade like positioning is modestly short \& in flattener; residual out of the money longs continue to bias rates to selloff}
\end{figure}
\begin{figure}[htbp]
\centering
\includegraphics[width=0.88\linewidth]{figures/fig_006.jpg}
\caption{Exhibit 2 - Exhibit 2: 2y contract CTA positioning across different lookback models All lookback models suggest CTAs are short 2y}
\end{figure}
\begin{figure}[htbp]
\centering
\includegraphics[width=0.88\linewidth]{figures/fig_007.jpg}
\caption{Exhibit 3 - Exhibit 3: 30y contract CTA positioning across different lookback models All lookback models suggest CTAs are short 30y}
\end{figure}
\begin{figure}[htbp]
\centering
\includegraphics[width=0.88\linewidth]{figures/fig_010.jpg}
\caption{Exhibit 7 - Exhibit 7: Asset manager and leveraged fund positioning (10y equivalent, \textbackslash{}\$bn) Asset manager longs correspond with leveraged fund shorts}
\end{figure}
\begin{figure}[htbp]
\centering
\includegraphics[width=0.88\linewidth]{figures/fig_013.jpg}
\caption{Exhibit 11 - Exhibit 11: Proxies for futures positioning (\textbackslash{}\$mil '01, duration-weighted by contract) In the money positions are predominantly shorts with some long out of the money in TY and WN points of the curve}
\end{figure}
\begin{figure}[htbp]
\centering
\includegraphics[width=0.88\linewidth]{figures/fig_020.jpg}
\caption{Exhibit 17 - Exhibit 17: UST supply versus sources of demand (\textbackslash{}\$bn) Demand in Q4 was weaker relative to Q3 despite similar supply}
\end{figure}
\begin{figure}[htbp]
\centering
\includegraphics[width=0.88\linewidth]{figures/fig_029.jpg}
\caption{Exhibit 27 - Exhibit 27: Japan investment in foreign bonds, cumulative weekly (\textbackslash{}\$bn) Long \& medium term bonds holdings rose \textbackslash{}\$10bn on the week}
\end{figure}
{\small\color{refgray}\textbf{References:} [R001] 市场真正在交易的不是AI，而是通胀预期重新定价; [R002] 市场对美联储缩表的担忧，可能高估了实际影响; [R003] 利率市场真正的拐点信号，不是降息时点，而是多头彻底离场}

\section{AI与算力基础设施}
\textbf{AI硬件供应链从需求侧转向供给侧，产能瓶颈和规格升级成为核心矛盾，光学引擎渗透率跃迁和CPU重新定价是最大预期差。}

\begin{itemize}
  \item TPU供应链扩容速度超预期，受益者从单一供应商向多层级扩散，CLS的TPU相关收入预计从2025年约10亿美元翻倍至2026年28亿美元。
  \item AI铜箔需求正从HVLP1/2/3向HVLP4加速迁移，单价可达传统铜箔10倍以上，供给缺口为中国供应商打开窗口期。
  \item NVL576架构使光学首次进入scale-up域，每GPU光学引擎附着率从2-4个跃升至17个，全光架构下将超过35个。
  \item Agentic CPU将从2025年约3亿美元飙升至2030年约594亿美元，年复合增长率185\%，到2030年将占据CPU总市场的45\%。
  \item 服务器CPU整体TAM预计从2025年293亿美元扩张至2030年1315亿美元，AMD被视为最大受益者。
  \item AI货币化正从概念验证进入规模兑现期，软件定价模式从按席位收费转向按Token收费，ERP和设计工具类增长可见度最高。
\end{itemize}
\begin{figure}[htbp]
\centering
\includegraphics[width=0.88\linewidth]{figures/fig_443.jpg}
\caption{Exhibit 1 - Exhibit 1: We expect the optical engines attach rate per GPU to increase 10x in future configurations}
\end{figure}
\begin{figure}[htbp]
\centering
\includegraphics[width=0.88\linewidth]{figures/fig_444.jpg}
\caption{Exhibit 2 - Exhibit 2: AI Transceivers global volume - Global market volume ('000 units)}
\end{figure}
\begin{figure}[htbp]
\centering
\includegraphics[width=0.88\linewidth]{figures/fig_448.jpg}
\caption{Exhibit 7 - Exhibit 7: Although transceivers are expected to show healthy growth... AI Transceivers global volume (m)}
\end{figure}
\begin{figure}[htbp]
\centering
\includegraphics[width=0.88\linewidth]{figures/fig_449.jpg}
\caption{Exhibit 8 - Exhibit 8: ...the larger opportunity lies in the adoption of CPO}
\end{figure}
\begin{figure}[htbp]
\centering
\includegraphics[width=0.88\linewidth]{figures/fig_451.jpg}
\caption{Exhibit 11 - Exhibit 11: Soitec Photonic SOI revenue to accelerate towards CY28e and maintain strong growth thereafter Soitec Photonics SOI revenue (\$m)}
\end{figure}
\begin{figure}[htbp]
\centering
\includegraphics[width=0.88\linewidth]{figures/fig_452.jpg}
\caption{Exhibit 12 - Exhibit 12: Besi hybrid bonding revenue Besi hybrid bonding revenue (€m)}
\end{figure}
{\small\color{refgray}\textbf{References:} [R006] 市场真正低估的不是AI需求，而是供给侧的约束正在重塑赢家; [R018] AI铜箔的真正瓶颈不是产能，而是规格升级的速度; [R029] AI 正在重塑中国软件业的定价权，但市场低估了“从 SaaS 到 RaaS”的深远影响; [R030] 市场真正低估的不是AI芯片，而是CPU的重新定价; [R036] 市场真正低估的不是光模块，而是光学引擎的渗透率跃迁}

\section{能源与大宗商品}
\textbf{供给侧不可逆重置正在重塑有色金属定价逻辑，铜铝市场结构性紧张远超市场定价，天然气和煤炭供给弹性收窄。}

\begin{itemize}
  \item 铜冶炼加工费TC/RC跌至-104美元/吨创历史新低，矿端议价权已彻底转移，冶炼厂倒贴钱买矿反映精矿供应极度紧张。
  \item 铝市场正经历近50年来最大供给冲击之一，超过300万吨年化产量损失，中国闲置产能已从35\%-47\%降至约5\%，供给弹性系统性下降。
  \item 即使需求疲软，未来6-12个月铝库存仍将刷新55年历史低点，2026下半年均价有望达4000美元/吨。
  \item 美国天然气Haynesville产区产量增长明显放缓，在均价3.19美元背景下产量反而下降，供给曲线可能比共识假设陡峭得多。
  \item 4月中国煤炭产量环比下降12.5\%且同比负增长，库存处于近年同期低位，夏季降温需求前供需缺口可能比预期更紧。
  \item 化工板块投资者仓位普遍偏低，市场正在定价2026上半年可能是未来两年利润高点，等待更清晰的定价锚。
\end{itemize}
\begin{figure}[htbp]
\centering
\includegraphics[width=0.88\linewidth]{figures/fig_127.jpg}
\caption{Exhibit 1 - Exhibit 1: Copper and aluminum spot prices LME copper price +2.8\% WoW to US\textbackslash{}\$13,895/t and LME aluminum price was +5.0\% WoW at USD 3,741/t as of May 15}
\end{figure}
\begin{figure}[htbp]
\centering
\includegraphics[width=0.88\linewidth]{figures/fig_128.jpg}
\caption{Exhibit 3 - Exhibit 3: Domestic monthly treatment charges on Cu concentrate Domestic monthly treatment charges on copper concentrate: -US\textbackslash{}\$103.72/t as of May 15}
\end{figure}
\begin{figure}[htbp]
\centering
\includegraphics[width=0.88\linewidth]{figures/fig_131.jpg}
\caption{Exhibit 4 - Exhibit 4: Avg. national aluminum margin (60\% captive plant) Avg. national aluminum margin rose by RMB 115/t WoW to RMB8,077/t as of May 15}
\end{figure}
\begin{figure}[htbp]
\centering
\includegraphics[width=0.88\linewidth]{figures/fig_244.jpg}
\caption{Figure 1 - The aluminium market is increasingly transitioning from an initial geopolitical shock into a structurally tighter inventory regime. The first phase of the move was dominated by war escalation risk, Strait of Hormuz concerns, oil price spikes, and broader macro}
\end{figure}
\begin{figure}[htbp]
\centering
\includegraphics[width=0.88\linewidth]{figures/fig_245.jpg}
\caption{Figure 2 - © 2026 Citi Inc. No redistribution without Citi's written permission. Figure 2. Historical spare capacity in Chian's aluminium industry has disappeared-implications for prices and supply elasticity}
\end{figure}
\begin{figure}[htbp]
\centering
\includegraphics[width=0.88\linewidth]{figures/fig_246.jpg}
\caption{Figure 11 - Only an extreme recession scenario comparable to the Volcker-era downturn or the 2008-09 GFC appears sufficient to stabilise aluminium inventory cover. We have already revised down our aluminium demand growth assumptions materially in our base case supply and }
\end{figure}
\begin{figure}[htbp]
\centering
\includegraphics[width=0.88\linewidth]{figures/fig_349.jpg}
\caption{Exhibit 2 - Exhibit 2: YTD growth in Haynesville production has been much lower than a year ago, suggesting a degree of sensitivity to price North Louisiana dry gas production, Bcf/d}
\end{figure}
\begin{figure}[htbp]
\centering
\includegraphics[width=0.88\linewidth]{figures/fig_386.jpg}
\caption{Exhibit 1 - Exhibit 1: Chart of the week – China: Monthly Coal Production}
\end{figure}
{\small\color{refgray}\textbf{References:} [R008] 铜冶炼加工费跌穿-100美元，市场真正低估的不是需求，而是供给侧的再定价; [R011] 铝市场最被低估的不是需求，而是供给侧的不可逆重置; [R014] 市场真正低估的不是需求，而是供给的价格弹性正在重塑天然气定价; [R015] 化工板块的真正考验不在油价，而在利润峰值能否被消化; [R019] 夏季用电高峰前，煤炭市场真正被低估的不是需求，而是供给侧的季节性收缩力度}

\section{地产与消费}
\textbf{中国房地产供给端正在发生结构性收缩，广义库存下降19\%，二手房市场活跃度远超新房；消费端补贴效应退潮后结构性分化加剧。}

\begin{itemize}
  \item 广义库存从2023年37.22亿平方米降至2025年30.2亿平方米，降幅约19\%，驱动因素是新开工持续萎缩而非销售加速。
  \item 4月住宅销售额同比降幅收窄至6.4\%，但环比下降33.7\%，改善几乎完全来自低基数效应而非需求实质性回暖。
  \item 二手房市场弹性远超新房，10城二手房周度成交量同比飙升49\%，一线城市同比大增65\%，定价权正从开发商向存量房业主转移。
  \item 4月社零同比增长仅0.2\%，零售商品在疫情后首次出现同比负增长，补贴效应退潮是核心原因，家电家居品类跌幅最大。
  \item 中国奶粉市场消费者决策从品牌出身转向成分科学，63\%因召回事件流失的进口品牌消费者转向国产品牌，55\%受访者对国内外品牌无偏好。
  \item 保险行业一季度盈利下滑几乎完全由投资端拖累，保险业务端实际稳定，寿险价值率在改善，市场对悲观已充分定价。
\end{itemize}
\begin{figure}[htbp]
\centering
\includegraphics[width=0.88\linewidth]{figures/fig_228.jpg}
\caption{Figure 3 - 4M26 REI -13.7\% ('25: -17.2\%); 4M26 FAI -1.6\% (NBS) Figure 3. Real Estate Investment: Monthly YoY Change (Apr-26)}
\end{figure}
\begin{figure}[htbp]
\centering
\includegraphics[width=0.88\linewidth]{figures/fig_234.jpg}
\caption{Figure 11 - © 2026 Citi Inc. No redistribution without Citi's written permission. Figure 11. National Residential Inventory hits 421m sqm by Apr 2026 (mn sqm GFA)}
\end{figure}
\begin{figure}[htbp]
\centering
\includegraphics[width=0.88\linewidth]{figures/fig_350.jpg}
\caption{Figure 5 - Figure 1: Broad inventory declines 19\% in 2024-25}
\end{figure}
\begin{figure}[htbp]
\centering
\includegraphics[width=0.88\linewidth]{figures/fig_351.jpg}
\caption{Figure 2 - Figure 2: Accumulated GFA started but not sold declined to 1.9bn sqm as of end-2025 from 2.5bn sqm as of end-2023}
\end{figure}
\begin{figure}[htbp]
\centering
\includegraphics[width=0.88\linewidth]{figures/fig_352.jpg}
\caption{Figure 3 - Figure 3: 19 cities had a \$10\textbackslash{}\%\$ -plus inventory decline in the past 12 months}
\end{figure}
\begin{figure}[htbp]
\centering
\includegraphics[width=0.88\linewidth]{figures/fig_398.jpg}
\caption{Exhibit 1 - Exhibit 1: Primary weekly unit sales and four-week moving average - Total 50 cities}
\end{figure}
\begin{figure}[htbp]
\centering
\includegraphics[width=0.88\linewidth]{figures/fig_387.jpg}
\caption{Exhibit 2 - Exhibit 3: MSCI China vs. Retail Sales Exhibit 4: China Retail Sales by Category Exhibit 6: Online Retail Sales}
\end{figure}
\begin{figure}[htbp]
\centering
\includegraphics[width=0.88\linewidth]{figures/fig_389.jpg}
\caption{Exhibit 4 - Exhibit 4: China Retail Sales by Category Exhibit 6: Online Retail Sales}
\end{figure}
{\small\color{refgray}\textbf{References:} [R009] 中国奶粉市场的真正拐点：消费者不再为“进口”买单，而是为“科学”付费; [R010] 中国房地产的真正拐点，不是销量，是库存结构; [R016] 中国房地产的真正拐点不在销售端，而在供给端; [R020] 市场低估的不是销售反弹，而是反弹的不可持续; [R021] 4月社零数据真正值得关注的不是0.2\%，而是“补贴效应退潮”后的结构性分化; [R022] 中国房地产的真正拐点，可能已经出现在二手房市场; [R025] 市场对保险板块的悲观已充分定价，但真正的上行风险来自结构而非弹性}

\section{区域市场与汽车}
\textbf{亚洲资本支出超级周期启动，日本经济缓冲垫耗尽，中国汽车出口成为利润锚点，全球电动车市场结构性分化加剧。}

\begin{itemize}
  \item 亚洲2026年GDP增速预期上调至4.8\%，工业周期超越能源冲击成为增长主导变量，四大结构性需求驱动资本开支超级周期。
  \item 日本Q1 GDP环比年化增长2.1\%超预期，但支撑经济韧性的缓冲垫正在消失，若油价进一步上行将暴露结构性脆弱性。
  \item 中国汽车4月国内零售同比下降21\%，出口却同比增长69\%且电动车占比提升至49\%，出口成为支撑产能利用率和盈利能力的关键。
  \item 全球电动车市场3月销量同比增长8\%但Q1累计仅增2.3\%，美国是唯一负增长主要市场（-35.5\%），欧洲反弹强劲。
  \item 全球汽车投资者情绪谨慎中带有明显分化，资金集中押注全球化执行能力、自动驾驶与物理AI、电池技术差异化三条叙事。
  \item 印度柴油汽油涨价不是终点而是亏损修复周期起点，运输燃料亏损已从15.5美元/桶收窄至9美元/桶，HPCL受益最大。
\end{itemize}
\begin{figure}[htbp]
\centering
\includegraphics[width=0.88\linewidth]{figures/fig_067.jpg}
\caption{Figure 1 - Figure 1: Real/nominal GDP (level)}
\end{figure}
\begin{figure}[htbp]
\centering
\includegraphics[width=0.88\linewidth]{figures/fig_070.jpg}
\caption{Figure 5 - Figure 5: Quarterly contribution (q/q, saar)}
\end{figure}
\begin{figure}[htbp]
\centering
\includegraphics[width=0.88\linewidth]{figures/fig_083.jpg}
\caption{Figure 18 - Figure 18: Nominal households consumption, disposable income and saving rate}
\end{figure}
\begin{figure}[htbp]
\centering
\includegraphics[width=0.88\linewidth]{figures/fig_086.jpg}
\caption{Figure 21 - Figure 21: Term of trade starts to worsen}
\end{figure}
\begin{figure}[htbp]
\centering
\includegraphics[width=0.88\linewidth]{figures/fig_088.jpg}
\caption{Figure 24 - Figure 24: The savings rate, which was relatively high in 2022, has recently turned negative, leaving limited room for further decline to support consumption}
\end{figure}
\begin{figure}[htbp]
\centering
\includegraphics[width=0.88\linewidth]{figures/fig_363.jpg}
\caption{Exhibit 5 - Long tail: 138 brands are competing for the remaining 38\% of the market (vs. 145 brands in 2025, 150 brands in 2024) \# Most frequently discussed charts of the month Exhibit 5. China passenger car export volume booked \$69\textbackslash{}\%\$ y-o-y growth in 4M26 China PC export}
\end{figure}
\begin{figure}[htbp]
\centering
\includegraphics[width=0.88\linewidth]{figures/fig_364.jpg}
\caption{Exhibit 6 - Exhibit 6. Chery, BYD and Geely lead China's passenger car export volume in 4M26}
\end{figure}
\begin{figure}[htbp]
\centering
\includegraphics[width=0.88\linewidth]{figures/fig_365.jpg}
\caption{Exhibit 7 - Exhibit 7. EV penetration rate rose to 61\% in April 2026}
\end{figure}
\begin{figure}[htbp]
\centering
\includegraphics[width=0.88\linewidth]{figures/fig_399.jpg}
\caption{Exhibit 1 - Exhibit 1: Fuel Prices in Asia and Refinery implications Exhibit 2: India Diesel Price Build-up after the fuel price hikes Price Build up - Diesel}
\end{figure}
{\small\color{refgray}\textbf{References:} [R004] 亚洲正在经历的，不是周期复苏，而是一个新的资本支出超级周期; [R005] 日本经济真正的拐点，不是GDP数字，而是“缓冲垫”已经耗尽; [R013] 全球电动车市场正在经历一场“结构性分化”，而非简单的复苏; [R017] 出口才是中国汽车利润的真正锚点; [R024] 印度燃料涨价的真正信号：零售商的亏损修复才刚刚开始; [R032] 全球汽车投资者真正焦虑的不是销量，而是“中国车企的下一步”还能被定价吗}

\section{风险偏好与资金流向}
\textbf{FOMO压倒宏观风险，利率与权益市场信号罕见背离，机构资金从押注AI赢家转向押注瓶颈环节。}

\begin{itemize}
  \item 利率波动率市场呈现陡峭看跌偏斜而权益看涨偏斜处于历史平坦水平，两个市场对同一宏观环境解读完全相反。
  \item 在泡沫构建阶段FOMO可以压倒收益率上行和宏观逆风，当前全球泡沫风险指标持续走高但尚未达到互联网泡沫峰值。
  \item 机构对Mega-cap科技股的低配幅度从-137个基点收窄至-125个基点，NVDA仍是最被低配个股（-2.39\%）。
  \item 存储与半导体设备公司SNDK、STX、LRCX、KLAC处于明显超配状态，机构资金正从押注最大市值AI赢家转向重仓基础设施瓶颈环节。
  \item 记忆体行业长期协议从软性承诺转向带有财务担保的硬约束，3-5年不可取消协议正在改写周期定价逻辑。
  \item 欧洲中小盘股价值重估不取决于宏观数据回暖，而取决于企业在细分赛道上的议价能力和战略执行力。
\end{itemize}
\begin{figure}[htbp]
\centering
\includegraphics[width=0.88\linewidth]{figures/fig_089.jpg}
\caption{Exhibit 7 - Exhibit 2: Latest\textbackslash{}* stress across GFSI sub-components Crude implied vol is the most stressed while Govt-OIS EUR is the least stressed}
\end{figure}
\begin{figure}[htbp]
\centering
\includegraphics[width=0.88\linewidth]{figures/fig_095.jpg}
\caption{Exhibit 8 - Exhibit 8: Bubble-like instability in global equities (MSCI World) broadly continues to rise, led by the BRIs of Kospi, Nikkei and US tech (though Mag7 still appears subdued); persistent geopolitical stress understandably keeps the}
\end{figure}
\begin{figure}[htbp]
\centering
\includegraphics[width=0.88\linewidth]{figures/fig_098.jpg}
\caption{Exhibit 8 - Exhibit 11: While the number of stocks exhibiting frothy price action has risen meaningfully, it remains localized vs Dotcom bubble peaks \# of SPX stocks with BRIs above 0.8}
\end{figure}
\begin{figure}[htbp]
\centering
\includegraphics[width=0.88\linewidth]{figures/fig_101.jpg}
\caption{Exhibit 15 - Exhibit 15: Unlike TLT skew, which has reached extreme steepness towards the put side, US equity vol skews are broadly looking through macro concerns, pricing in resilience and greater call side risks in the near-term 1m 25d/ATMf pu}
\end{figure}
\begin{figure}[htbp]
\centering
\includegraphics[width=0.88\linewidth]{figures/fig_102.jpg}
\caption{Exhibit 16 - Exhibit 16: Bubble-like regimes can withstand rising yields, as seen from the strong Nasdaq rally in the latter stages of the Dotcom bubble alongside higher 30y yields NDX \& US 30y yields during the 1990s Dotcom bubble}
\end{figure}
\begin{figure}[htbp]
\centering
\includegraphics[width=0.88\linewidth]{figures/fig_429.jpg}
\caption{Exhibit1 - Exhibit1: The gap between mega-cap tech's institutional weighting and its S\&P 500 weighting narrowed in 1Q26. Average Institutional Portfolio Weighting}
\end{figure}
\begin{figure}[htbp]
\centering
\includegraphics[width=0.88\linewidth]{figures/fig_430.jpg}
\caption{Exhibit 2 - Exhibit 2: Of the large cap stocks we evaluate, NVDA, AAPL, MSFT, AMZN, and GOOGL are currently the most under-owned in actively managed portfolios vs. the S\&P 500, while SNDK, STX, CRM, and LRCX are the most over-owned. Average Ac}
\end{figure}
\begin{figure}[htbp]
\centering
\includegraphics[width=0.88\linewidth]{figures/fig_431.jpg}
\caption{Exhibit 3 - Exhibit 3: Of the large cap stocks we evaluate, SNDK saw the gap between its institutional ownership and S\&P 500 weighting increase the most by 57bps in 1Q (vs. 4Q), while KLAC saw the gap narrow by 30bps. Q/Q Change in Gap of Acti}
\end{figure}
{\small\color{refgray}\textbf{References:} [R007] 市场真正低估的不是需求，而是供给侧的再定价; [R031] 机构持仓数据揭示的真正信号：AI投资从“押注赢家”转向“押注瓶颈”; [R034] 记忆体行业的“结构性契约”正在改写周期的定价逻辑; [R035] 欧洲中小盘股的真正机会，藏在管理层问答的缝隙里}

\section*{结语}
当前市场核心矛盾在于通胀预期重夺定价权与AI叙事泡沫化之间的碰撞，供给侧约束成为多板块共同主线。投资者需关注利率市场多头叙事终结后的定价逻辑切换，以及亚洲资本支出超级周期、有色金属供给不可逆重置、中国房地产库存结构改善等结构性变化带来的中长期影响。
\vfill
{\small\color{refgray}Personal reading notes and learning share only. Not investment advice.}
\end{document}
